% =====================================================================
\section{Offline}
\begin{frame}[plain]
  \centering
  \vfill
  {\Large\bfseries\color{accent} Part I --- Offline PAC Labeling}\\[6pt]
  \normalsize Certify one price on a calibration set, then label the whole pool.
  \vfill
\end{frame}

% ----------------------------------------------- OFFLINE CALIBRATION
\begin{frame}{Offline: certify a price from a small labeled sample}
  \small
  \textbf{Fixed-price rule.} If true losses were known, then once
  \(\lambda\) is fixed, the decision separates across items:
  \[
      a_t^\lambda
      =
      \argmin_{a\in\calA}
      \{c_t(a)+\lambda h_t(a)\}.
  \]
  \begin{itemize}
    \item But accessing \(h_t(a)\) requires the true label \(Y_t\).
    \item So we form a plug-in policy family using fixed proxy losses:
      \[
          a_t^\lambda
          =
          \argmin_{a\in\calA}
          \{c_t(a)+\lambda \hath_t(a)\}.
      \]
    \item The remaining question is how to choose \(\lambda\).
  \end{itemize}
  \begin{block}{Calibration idea}
    Label a small random subset from the whole pool and choose
    \(\hat\lambda\) to minimize cost while keeping error \(\le\eps\) with high
    probability.
  \end{block}
\end{frame}

% ----------------------------------------------- OFFLINE CALIBRATION
\begin{frame}{Offline: compute the certified price}
  \small
  \setlength{\abovedisplayskip}{4pt}
  \setlength{\belowdisplayskip}{4pt}
  \begin{enumerate}
    \item Draw calibration indices \(I_1,\ldots,I_m\) from the full pool and
          query their labels.
    \item Estimate the error of every candidate price:
      \[
          \hatL_m(\lambda)=\frac1m\sum_{s=1}^m h_{I_s}(a_{I_s}^\lambda),
          \qquad
          U_m(\lambda;\alpha)=
          \hatL_m(\lambda)+\sqrt{\frac{\log(1/\alpha)}{2m}}.
      \]
    \item Choose the first price whose upper bound is below the target:
      \[
          \hat\lambda
          =
          \min\{\lambda\ge0:\; U_m(\lambda;\alpha)\le\eps\}.
      \]
  \end{enumerate}
  \vspace{-1mm}
  \begin{block}{Why this is enough}
    In the equal-cost rule, increasing \(\lambda\) only moves items
    \emph{AI \(\to\) expert}, so the policies are nested and
    \(L_T(\lambda)\) is monotone. The next slide states the resulting PAC
    guarantee for the released labels.
  \end{block}
\end{frame}

% --------------------------------------------------- OFFLINE GUARANTEE
\begin{frame}{Offline: the PAC guarantee}
  \begin{block}{Theorem (offline monotone-threshold PAC)}
    With probability at least $1-\alpha$ over the calibration draw, the released
    labels satisfy
    \[
       \frac1T\sum_{t=1}^{T}\ell(Y_t,\tildeY_t)\;\le\;\eps.
    \]
  \end{block}
  \begin{itemize}
    \item Distribution-free, finite-sample. Falls back to expert-only if no
          finite price qualifies.
    \item Proxies affect \emph{cost}, not \emph{validity}: a bad proxy just
          defers more often.
  \end{itemize}
\end{frame}

% ------------------------------------------------ OFFLINE EXPERIMENT 1
\begin{frame}{Offline experiments: LP rule matches original PAC}
  \small
  \textbf{One cheap source} (\(\eps=0.10\)):
  \begin{table}
  \centering
  \begin{tabular}{llcc}
    \toprule
    Dataset & Method & Error & Budget saved \\
    \midrule
    ImageNet   & LP fixed price & 0.087 & 0.785 \\
    ImageNet   & Original PAC   & 0.088 & 0.786 \\
    ImageNetV2 & LP fixed price & 0.072 & 0.546 \\
    ImageNetV2 & Original PAC   & 0.073 & 0.548 \\
    \bottomrule
  \end{tabular}
  \end{table}
  \begin{itemize}
    \item Error stays under target \(\eps\) in all reported runs.
    \item With one cheap source, the LP fixed-price rule recovers the original
          PAC threshold behavior.
    \item On ImageNet, both methods save about \(78\%\) of expert budget.
  \end{itemize}
\end{frame}

% ------------------------------------------------ OFFLINE EXPERIMENT 2
\begin{frame}{Offline experiments: routing beats single sources}
  \small
  \textbf{Two cheap sources} (\(\eps=0.10\)):
  \vspace{-1mm}
  \begin{center}
  \scriptsize
  \begin{tabular}{llcc}
    \toprule
    Dataset & Method & Error & Budget saved \\
    \midrule
    ImageNet   & LP routing          & 0.087 & \textbf{0.785} \\
    ImageNet   & Best single source  & 0.088 & 0.475 \\
    ImageNetV2 & LP routing          & 0.072 & \textbf{0.546} \\
    ImageNetV2 & Best single source  & 0.072 & 0.359 \\
    \bottomrule
  \end{tabular}
  \end{center}
  \vspace{-2mm}
  \begin{center}
  \begin{tikzpicture}[
      scale=0.9,
      transform shape,
      font=\tiny,
      >=Latex,
      every node/.style={font=\tiny},
      card/.style={draw=accent!55, line width=0.6pt, rounded corners=3pt,
                   fill=accent!4, align=center, inner sep=0pt},
      title/.style={fill=accent, text=white, font=\tiny\bfseries,
                    rounded corners=2pt, inner xsep=5pt, inner ysep=2pt},
      arrow/.style={-{Latex[length=1.8mm,width=1.6mm]}, line width=0.9pt,
                    draw=accent!75},
      ev/.style={text=evenc, font=\tiny\bfseries},
      od/.style={text=oddc,  font=\tiny\bfseries}
    ]
    % ---- ORIGINAL SINGLE SOURCE --------------------------------------
    \node[card, minimum width=3.7cm] (single) at (0,0) {
      \setlength{\tabcolsep}{4pt}\renewcommand{\arraystretch}{1.15}%
      \begin{tabular}{@{}lll@{}}
        {\color{accent!60!black}Group} & {\color{accent!60!black}Label} &
          {\color{accent!60!black}Confidence} \\[0.3mm]
        \arrayrulecolor{accent!30}\hline\\[-2.2mm]
        {\color{evenc}\textbf{Even}} & \kept{orig.\ \(\hatY_t\)} & \confhi\,orig. \\
        {\color{oddc}\textbf{Odd}}   & \kept{orig.\ \(\hatY_t\)} & \confhi\,orig.
      \end{tabular}
    };
    \node[title, anchor=south] at (single.north) {Original single source};

    % ---- SOURCE A -----------------------------------------------------
    \node[card, minimum width=3.7cm] (sourceA) at (6.0,1.15) {
      \setlength{\tabcolsep}{4pt}\renewcommand{\arraystretch}{1.15}%
      \begin{tabular}{@{}lll@{}}
        {\color{accent!60!black}Group} & {\color{accent!60!black}Label} &
          {\color{accent!60!black}Confidence} \\[0.3mm]
        \arrayrulecolor{accent!30}\hline\\[-2.2mm]
        {\color{evenc}\textbf{Even}} & \kept{orig.\ \(\hatY_t\)} & \confhi\,orig. \\
        {\color{oddc}\textbf{Odd}}   & \rand{random}            & \conflo\,low
      \end{tabular}
    };
    \node[title, anchor=south] at (sourceA.north) {Source A \;\textcolor{white!55}{$\vert$}\; keeps \textcolor{evenc!30!white}{Even}};

    % ---- SOURCE B -----------------------------------------------------
    \node[card, minimum width=3.7cm] (sourceB) at (6.0,-1.15) {
      \setlength{\tabcolsep}{4pt}\renewcommand{\arraystretch}{1.15}%
      \begin{tabular}{@{}lll@{}}
        {\color{accent!60!black}Group} & {\color{accent!60!black}Label} &
          {\color{accent!60!black}Confidence} \\[0.3mm]
        \arrayrulecolor{accent!30}\hline\\[-2.2mm]
        {\color{evenc}\textbf{Even}} & \rand{random}            & \conflo\,low \\
        {\color{oddc}\textbf{Odd}}   & \kept{orig.\ \(\hatY_t\)} & \confhi\,orig.
      \end{tabular}
    };
    \node[title, anchor=south] at (sourceB.north) {Source B \;\textcolor{white!55}{$\vert$}\; keeps \textcolor{oddc!55!white}{Odd}};

    % ---- branching arrows --------------------------------------------
    \draw[arrow] (single.east) to[out=15,in=180]  (sourceA.west);
    \draw[arrow] (single.east) to[out=-15,in=180] (sourceB.west);
    \node[fill=accent!10, draw=accent!40, rounded corners=2pt,
          inner sep=2.5pt, font=\tiny\itshape] at (3.05,0)
      {split by true-class parity};
	  \end{tikzpicture}
	  \vspace{-1mm}

	  {\tiny \(\textnormal{random}\sim\mathrm{Uniform}(0,1)\), \quad
	  \(\textnormal{low confidence}\sim\mathrm{Uniform}(0,0.2)\).}
	  \end{center}
\end{frame}
